\documentclass[12pt,a4paper]{article}

% ------------------------------------------------
% Paketler
% ------------------------------------------------
\usepackage[utf8]{inputenc}
\usepackage[T1]{fontenc}

\usepackage{geometry}
\usepackage{array}
\usepackage{tabularx}
\usepackage{longtable}
\usepackage{booktabs}
\usepackage{float}
\usepackage{pdflscape}
\usepackage{ragged2e}
\usepackage{xcolor}
\usepackage{hyperref}

% ------------------------------------------------
% Sayfa ayarları
% ------------------------------------------------
\geometry{
    left=2.5cm,
    right=2.5cm,
    top=2.5cm,
    bottom=2.5cm
}

\setlength{\parindent}{0pt}
\setlength{\parskip}{6pt}

% ------------------------------------------------
% Tablo sütun tipleri
% ------------------------------------------------
\newcolumntype{L}[1]{>{\RaggedRight\arraybackslash}p{#1}}
\newcolumntype{C}[1]{>{\Centering\arraybackslash}p{#1}}
\newcolumntype{Y}{>{\RaggedRight\arraybackslash}X}

% ------------------------------------------------
% Hyperlink ayarları
% ------------------------------------------------
\hypersetup{
    colorlinks=true,
    linkcolor=black,
    urlcolor=blue
}

% ------------------------------------------------
% Başlık
% ------------------------------------------------
\title{
    \textbf{DEOS ROS 2 Topic Arayüz Tasarımı}
}
\author{\textbf{Emir Arkalı} (\href{mailto:emirarkali07@gmail.com}{emirarkali07@gmail.com})\\DEOS -- Dokuz Eylül Otonom Sistemler}
\date{İlk Yayın: 28 Eylül 2026 \quad | \quad Son Revizyon: 13 Ekim 2026}

% ================================================================
\begin{document}
% ================================================================

\maketitle

\section{Amaç}

Bu dokümanın amacı, araç üzerindeki ROS 2 yazılım bileşenleri ile
\texttt{deos\_ros\_bridge} düğümü arasındaki topic tabanlı haberleşme
arayüzünü tanımlamaktır.

\texttt{deos\_ros\_bridge}, ROS 2 tarafı ile araç üzerindeki gömülü kontrol
altyapısı arasında geçiş katmanı olarak görev yapacaktır.

ROS 2 açısından bridge iki temel göreve sahiptir:

\begin{itemize}
    \item ROS 2 içerisindeki kontrol ve otonom sürüş düğümleri tarafından
    publish edilen komut topic'lerine subscribe olmak,
    
    \item Araçtan alınan durum ve sensör bilgilerini ilgili ROS 2 topic'leri
    üzerinden publish etmek.
\end{itemize}

ROS 2 tarafında fiziksel büyüklükler SI birimleri kullanılarak taşınacaktır.
ROS 2 düğümlerinde fixed-point ölçekleme kullanılmayacaktır.

Gerekli ölçekleme veya gömülü sistem veri formatına dönüşüm,
\texttt{deos\_ros\_bridge} içerisinde gerçekleştirilecektir.

% ================================================================
\section{ROS 2'den Bridge'e Gönderilen Topic'ler}
% ================================================================

Aşağıdaki topic'ler ROS 2 tarafındaki otonom sürüş veya kontrol düğümleri
tarafından publish edilecektir.

\texttt{deos\_ros\_bridge} bu topic'lere subscriber olarak bağlanacaktır.

% ================================================================
% LANDSCAPE TABLO 1
% ================================================================

\clearpage

\begin{landscape}

\thispagestyle{plain}

\begin{table}[p]

\centering

\small

\renewcommand{\arraystretch}{1.35}

\begin{tabularx}{\linewidth}{
|L{5.2cm}
|L{5.0cm}
|L{4.0cm}
|L{2.2cm}
|Y|
}

\hline

\textbf{Topic}
&
\textbf{Message Type}
&
\textbf{Birim / İçerik}
&
\textbf{Rate}
&
\textbf{QoS}
\\

\hline

\texttt{/cmd\_vel}
&
\texttt{geometry\_msgs/msg/Twist}
&
\texttt{linear.x}: m/s

\texttt{angular.z}: rad/s
&
20 Hz
&
RELIABLE

KEEP\_LAST(1)

VOLATILE
\\

\hline

\texttt{/vehicle/autonomy\_enable}
&
\texttt{std\_msgs/msg/Bool}
&
Boolean
&
5 Hz
&
RELIABLE

KEEP\_LAST(1)

VOLATILE
\\

\hline

\texttt{/vehicle/motion\_enable}
&
\texttt{std\_msgs/msg/Bool}
&
Boolean
&
10 Hz
&
RELIABLE

KEEP\_LAST(1)

VOLATILE
\\

\hline

\texttt{/vehicle/signals/right}
&
\texttt{std\_msgs/msg/Bool}
&
Boolean
&
On Change
&
RELIABLE

KEEP\_LAST(1)

TRANSIENT\_LOCAL
\\

\hline

\texttt{/vehicle/signals/left}
&
\texttt{std\_msgs/msg/Bool}
&
Boolean
&
On Change
&
RELIABLE

KEEP\_LAST(1)

TRANSIENT\_LOCAL
\\

\hline

\texttt{/vehicle/signals/hazard}
&
\texttt{std\_msgs/msg/Bool}
&
Boolean
&
On Change
&
RELIABLE

KEEP\_LAST(1)

TRANSIENT\_LOCAL
\\

\hline

\texttt{/vehicle/lights/brake}
&
\texttt{std\_msgs/msg/Bool}
&
Boolean
&
On Change
&
RELIABLE

KEEP\_LAST(1)

TRANSIENT\_LOCAL
\\

\hline

\texttt{/vehicle/lights/headlight}
&
\texttt{std\_msgs/msg/Bool}
&
Boolean
&
On Change
&
RELIABLE

KEEP\_LAST(1)

TRANSIENT\_LOCAL
\\

\hline

\end{tabularx}

\caption{ROS 2 tarafından publish edilip bridge tarafından subscribe edilen topic'ler}

\end{table}

\end{landscape}

\clearpage

% ================================================================
\subsection{\texttt{/cmd\_vel}}
% ================================================================

Bu topic, aracın otonom sürüş sırasında uygulaması gereken hareket komutunu
taşımaktadır.

Mesaj tipi:

\begin{verbatim}
geometry_msgs/msg/Twist
\end{verbatim}

Mesaj içerisindeki aşağıdaki alanlar kullanılacaktır:

\begin{itemize}
    \item \texttt{linear.x}: hedef doğrusal hız,
    \item \texttt{angular.z}: hedef yaw açısal hızı.
\end{itemize}

Birimler:

\begin{verbatim}
linear.x  -> m/s
angular.z -> rad/s
\end{verbatim}

Diğer alanlar kullanılmayacak ve sıfır değerinde tutulacaktır:

\begin{verbatim}
linear.y  = 0
linear.z  = 0

angular.x = 0
angular.y = 0
\end{verbatim}

Topic yayın frekansı:

\begin{verbatim}
20 Hz
\end{verbatim}

QoS profili:

\begin{verbatim}
Reliability : RELIABLE
History     : KEEP_LAST
Depth       : 1
Durability  : VOLATILE
\end{verbatim}

Bridge tarafında son geçerli \texttt{/cmd\_vel} mesajı için timeout
mekanizması uygulanacaktır.

Önerilen timeout değeri:

\begin{verbatim}
200 ms
\end{verbatim}

Bu süre içerisinde yeni bir \texttt{/cmd\_vel} mesajı alınmazsa,
önceki hareket komutu geçersiz kabul edilmelidir.

% ================================================================
\subsection{\texttt{/vehicle/autonomy\_enable}}
% ================================================================

Bu topic, otonom sürüş fonksiyonunun aktif tutulmasının istenip
istenmediğini belirtmektedir.

Mesaj tipi:

\begin{verbatim}
std_msgs/msg/Bool
\end{verbatim}

Anlamı:

\begin{verbatim}
true  -> Otonom sürüş aktif
false -> Otonom sürüş pasif
\end{verbatim}

Yayın frekansı:

\begin{verbatim}
5 Hz
\end{verbatim}

QoS profili:

\begin{verbatim}
Reliability : RELIABLE
History     : KEEP_LAST
Depth       : 1
Durability  : VOLATILE
\end{verbatim}

Önerilen timeout değeri:

\begin{verbatim}
500 ms
\end{verbatim}

Belirlenen süre içerisinde yeni mesaj alınmaması durumunda:

\begin{verbatim}
autonomy_enable = false
\end{verbatim}

olarak değerlendirilmelidir.

% ================================================================
\subsection{\texttt{/vehicle/motion\_enable}}
% ================================================================

Bu topic, aracın fiziksel olarak hareket etmesine izin verilip
verilmediğini belirtmektedir.

Mesaj tipi:

\begin{verbatim}
std_msgs/msg/Bool
\end{verbatim}

Anlamı:

\begin{verbatim}
true  -> Araç hareket komutlarını uygulayabilir
false -> Araç hareket ettirilmemelidir
\end{verbatim}

Yayın frekansı:

\begin{verbatim}
10 Hz
\end{verbatim}

QoS profili:

\begin{verbatim}
Reliability : RELIABLE
History     : KEEP_LAST
Depth       : 1
Durability  : VOLATILE
\end{verbatim}

Önerilen timeout:

\begin{verbatim}
250 ms
\end{verbatim}

Belirlenen süre içerisinde yeni mesaj alınmazsa:

\begin{verbatim}
motion_enable = false
\end{verbatim}

olarak değerlendirilmelidir.

Bu davranış, haberleşme veya ROS 2 yazılım hatası durumunda aracın
hareket komutlarının süresiz olarak uygulanmasını engellemek amacıyla
kullanılacaktır.

% ================================================================
\section{Sinyal ve Aydınlatma Topic'leri}
% ================================================================

Araç üzerindeki sinyal ve aydınlatma sistemlerini kontrol etmek amacıyla
aşağıdaki topic'ler kullanılacaktır:

\begin{itemize}

    \item \texttt{/vehicle/signals/right}

    \item \texttt{/vehicle/signals/left}

    \item \texttt{/vehicle/signals/hazard}

    \item \texttt{/vehicle/lights/brake}

    \item \texttt{/vehicle/lights/headlight}

\end{itemize}

Bu topic'lerin tamamı aşağıdaki mesaj tipini kullanacaktır:

\begin{verbatim}
std_msgs/msg/Bool
\end{verbatim}

Boolean değerlerin anlamı:

\begin{verbatim}
true  -> İlgili çıkış aktif
false -> İlgili çıkış pasif
\end{verbatim}

Bu topic'ler yalnızca durum değiştiğinde publish edilebilir.

Yayın davranışı:

\begin{verbatim}
On Change
\end{verbatim}

QoS profili:

\begin{verbatim}
Reliability : RELIABLE
History     : KEEP_LAST
Depth       : 1
Durability  : TRANSIENT_LOCAL
\end{verbatim}

\texttt{TRANSIENT\_LOCAL} kullanılması sayesinde ROS 2 ağına sonradan
bağlanan bir subscriber ilgili sinyal veya lambanın son publish edilen
durumunu alabilecektir.

% ================================================================
\section{Bridge'den ROS 2'ye Gönderilen Topic'ler}
% ================================================================

Araç üzerindeki gömülü kontrol sisteminden alınan durum ve sensör bilgileri
\texttt{deos\_ros\_bridge} tarafından aşağıdaki ROS 2 topic'leri üzerinden
publish edilecektir.

% ================================================================
% LANDSCAPE TABLO 2
% ================================================================

\clearpage

\begin{landscape}

\thispagestyle{plain}

\begin{table}[p]

\centering

\small

\renewcommand{\arraystretch}{1.4}

\begin{tabularx}{\linewidth}{
|L{5.5cm}
|L{5.2cm}
|L{3.2cm}
|L{2.3cm}
|Y|
}

\hline

\textbf{Topic}
&
\textbf{Message Type}
&
\textbf{Birim}
&
\textbf{Rate}
&
\textbf{QoS}
\\

\hline

\texttt{/vehicle/state/speed}
&
\texttt{std\_msgs/msg/Float32}
&
m/s
&
20 Hz
&
BEST\_EFFORT

KEEP\_LAST(1)

VOLATILE
\\

\hline

\texttt{/vehicle/state/steering\_angle}
&
\texttt{std\_msgs/msg/Float32}
&
rad
&
20 Hz
&
BEST\_EFFORT

KEEP\_LAST(1)

VOLATILE
\\

\hline

\texttt{/sensors/imu/data}
&
\texttt{sensor\_msgs/msg/Imu}
&
SI Units
&
50 Hz
&
BEST\_EFFORT

KEEP\_LAST(5)

VOLATILE
\\

\hline

\texttt{/sensors/distance/<sensor\_name>}
&
\texttt{sensor\_msgs/msg/Range}
&
m
&
20 Hz
&
BEST\_EFFORT

KEEP\_LAST(1)

VOLATILE
\\

\hline

\texttt{/vehicle/state\_machine/state}
&
\texttt{deos\_msgs/msg/VehicleState}
&
Enum
&
On Change
&
RELIABLE

KEEP\_LAST(1)

TRANSIENT\_LOCAL
\\

\hline

\end{tabularx}

\caption{Bridge tarafından ROS 2 sistemine publish edilen topic'ler}

\end{table}

\end{landscape}

\clearpage

% ================================================================
\section{Anlık Hız Bilgisi}
% ================================================================

Topic adı:

\begin{verbatim}
/vehicle/state/speed
\end{verbatim}

Mesaj tipi:

\begin{verbatim}
std_msgs/msg/Float32
\end{verbatim}

Payload:

\begin{verbatim}
float32 data
\end{verbatim}

Birim:

\begin{verbatim}
m/s
\end{verbatim}

İşaret konvansiyonu:

\begin{verbatim}
Pozitif değer -> İleri hareket
0             -> Araç duruyor
Negatif değer -> Geri hareket
\end{verbatim}

Publish frekansı:

\begin{verbatim}
20 Hz
\end{verbatim}

Bu frekans, her:

\begin{verbatim}
50 ms
\end{verbatim}

aralıkta yeni hız bilgisi yayınlanmasına karşılık gelmektedir.

QoS profili:

\begin{verbatim}
Reliability : BEST_EFFORT
History     : KEEP_LAST
Depth       : 1
Durability  : VOLATILE
\end{verbatim}

Bu topic için yalnızca en güncel hız değeri önemlidir.
Eski hız verilerinin tekrar iletilmesine ihtiyaç bulunmamaktadır.

% ================================================================
\section{Anlık Direksiyon Açısı}
% ================================================================

Topic adı:

\begin{verbatim}
/vehicle/state/steering_angle
\end{verbatim}

Mesaj tipi:

\begin{verbatim}
std_msgs/msg/Float32
\end{verbatim}

Birim:

\begin{verbatim}
rad
\end{verbatim}

Açı işaret konvansiyonu:

\begin{verbatim}
0        -> Direksiyon düz
Pozitif  -> Sola dönüş
Negatif  -> Sağa dönüş
\end{verbatim}

Örnek:

\begin{verbatim}
+0.2618 rad ~= +15 degree
-0.2618 rad ~= -15 degree
\end{verbatim}

Publish frekansı:

\begin{verbatim}
20 Hz
\end{verbatim}

QoS profili:

\begin{verbatim}
Reliability : BEST_EFFORT
History     : KEEP_LAST
Depth       : 1
Durability  : VOLATILE
\end{verbatim}

Bu topic içerisinde yalnızca direksiyon sisteminin en güncel fiziksel açı
bilgisinin korunması amaçlanmaktadır.

% ================================================================
\section{IMU Verisi}
% ================================================================

Topic adı:

\begin{verbatim}
/sensors/imu/data
\end{verbatim}

Mesaj tipi:

\begin{verbatim}
sensor_msgs/msg/Imu
\end{verbatim}

Mesajın temel alanları aşağıdaki gibidir:

\begin{itemize}

    \item Header

    \item Orientation quaternion

    \item Angular velocity

    \item Linear acceleration

    \item Orientation covariance

    \item Angular velocity covariance

    \item Linear acceleration covariance

\end{itemize}

Birimler:

\begin{verbatim}
Angular Velocity     -> rad/s

Linear Acceleration  -> m/s^2

Orientation          -> Quaternion
\end{verbatim}

Publish frekansı:

\begin{verbatim}
50 Hz
\end{verbatim}

Bu frekans:

\begin{verbatim}
20 ms
\end{verbatim}

örnekleme/yayın periyoduna karşılık gelmektedir.

QoS profili:

\begin{verbatim}
Reliability : BEST_EFFORT
History     : KEEP_LAST
Depth       : 5
Durability  : VOLATILE
\end{verbatim}

IMU verisinin sürekli ve yüksek frekanslı olması nedeniyle
\texttt{BEST\_EFFORT} profili tercih edilmiştir.

% ================================================================
\section{Mesafe Sensörleri}
% ================================================================

Mesafe sensörleri aşağıdaki namespace yapısını kullanacaktır:

\begin{verbatim}
/sensors/distance/<sensor_name>
\end{verbatim}

Buradaki:

\begin{verbatim}
<sensor_name>
\end{verbatim}

ifadesi sensörün araç üzerindeki fiziksel konumunu belirtecektir.

Örnek topic isimleri:

\begin{verbatim}
/sensors/distance/front

/sensors/distance/front_left

/sensors/distance/front_right

/sensors/distance/rear
\end{verbatim}

Sensörlerin kesin sayısı ve fiziksel yerleşimi henüz belirlenmediğinden
bu topic isimleri placeholder olarak tanımlanmıştır.

Sensör yerleşimi kesinleştirildiğinde her sensör için gerçek topic adı
belirlenecektir.

Mesaj tipi:

\begin{verbatim}
sensor_msgs/msg/Range
\end{verbatim}

Mesajın kullanılabilecek temel alanları:

\begin{verbatim}
radiation_type

field_of_view

min_range

max_range

range
\end{verbatim}

Mesafe birimi:

\begin{verbatim}
m
\end{verbatim}

Başlangıç publish frekansı:

\begin{verbatim}
20 Hz
\end{verbatim}

QoS profili:

\begin{verbatim}
Reliability : BEST_EFFORT
History     : KEEP_LAST
Depth       : 1
Durability  : VOLATILE
\end{verbatim}

Bu topic'lerde eski mesafe ölçümlerinin tekrar iletilmesi yerine
en güncel ölçüm değerinin kullanılmasına öncelik verilecektir.

% ================================================================
\section{Araç State Machine Bilgisi}
% ================================================================

Araç kontrol sisteminin mevcut çalışma durumu aşağıdaki topic üzerinden
ROS 2 ağına aktarılacaktır:

\begin{verbatim}
/vehicle/state_machine/state
\end{verbatim}

Mesaj tipi:

\begin{verbatim}
deos_msgs/msg/VehicleState
\end{verbatim}

Bu mesaj için önerilen yapı aşağıdaki gibidir:

\begin{verbatim}
std_msgs/Header header

uint8 UNKNOWN=0
uint8 INIT=1
uint8 IDLE=2
uint8 MANUAL=3
uint8 AUTONOMOUS=4
uint8 SAFETY=5
uint8 EMERGENCY=6
uint8 FAULT=7

uint8 state
\end{verbatim}

Mevcut state değerleri state machine mimarisi kesinleştirildiğinde
güncellenebilir.

Topic sürekli periyodik olarak yayınlanmayacaktır.

Publish davranışı:

\begin{verbatim}
On Change
\end{verbatim}

Örneğin araç:

\begin{verbatim}
AUTONOMOUS
\end{verbatim}

durumundan:

\begin{verbatim}
SAFETY
\end{verbatim}

durumuna geçtiğinde yeni state değeri publish edilecektir.

QoS profili:

\begin{verbatim}
Reliability : RELIABLE
History     : KEEP_LAST
Depth       : 1
Durability  : TRANSIENT_LOCAL
\end{verbatim}

\texttt{TRANSIENT\_LOCAL} kullanılması sayesinde ROS 2 ağına sonradan
bağlanan bir node, yeni bir state değişikliği gerçekleşmesini beklemeden
aracın son publish edilen state bilgisini alabilecektir.

% ================================================================
\section{Topic Yönlerinin Özeti}
% ================================================================

\subsection{Bridge Tarafından Subscribe Edilen Topic'ler}

ROS 2 kontrol ve otonom sürüş düğümleri tarafından publish edilmesi
gereken topic'ler:

\begin{verbatim}
/cmd_vel

/vehicle/autonomy_enable

/vehicle/motion_enable

/vehicle/signals/right

/vehicle/signals/left

/vehicle/signals/hazard

/vehicle/lights/brake

/vehicle/lights/headlight
\end{verbatim}

\subsection{Bridge Tarafından Publish Edilen Topic'ler}

Araçtan gelen durum ve sensör verilerinin yayınlanacağı topic'ler:

\begin{verbatim}
/vehicle/state/speed

/vehicle/state/steering_angle

/sensors/imu/data

/sensors/distance/<sensor_name>

/vehicle/state_machine/state
\end{verbatim}

% ================================================================
\section{Genel ROS 2 Veri Akışı}
% ================================================================

ROS 2 tarafındaki veri akışı genel olarak aşağıdaki yapıdadır:

\begin{verbatim}

             ROS 2 AUTONOMY / CONTROL

                      |
                      |
                   publish
                      |
                      v

              +----------------+
              |                |
              | deos_ros_bridge|
              |                |
              +----------------+

                      |
                      |
                   publish
                      |
                      v

          ROS 2 STATE / SENSOR CONSUMERS


\end{verbatim}

ROS 2 kontrol düğümleri araç donanımıyla doğrudan haberleşmeyecektir.

Araç ile ilgili komutların tamamı tanımlanan ROS 2 topic'leri üzerinden
\texttt{deos\_ros\_bridge} düğümüne aktarılacaktır.

Benzer şekilde araçtan alınan durum ve sensör bilgileri bridge tarafından
ROS 2 topic'lerine dönüştürülecektir.

% ================================================================
\section{QoS Terimleri}
% ================================================================

Bu dokümanda kullanılan temel ROS 2 QoS parametreleri aşağıda
özetlenmiştir.

\subsection{RELIABLE}

Mesajların subscriber'a ulaştırılmasına öncelik verilir.

Durum değişiklikleri, enable bilgileri ve state machine gibi kaybının
istenmediği mesajlarda kullanılmaktadır.

\subsection{BEST\_EFFORT}

Paket kaybı oluşması durumunda eski mesajın tekrar gönderilmesi yerine
yeni mesajların iletilmesine devam edilir.

Sürekli güncellenen hız, direksiyon açısı, IMU ve mesafe ölçümleri için
kullanılmaktadır.

\subsection{KEEP\_LAST}

Subscriber kuyruğunda yalnızca belirtilen sayıdaki en son mesaj tutulur.

Örneğin:

\begin{verbatim}
KEEP_LAST
Depth = 1
\end{verbatim}

kullanılması durumunda yalnızca son mesaj saklanır.

\subsection{VOLATILE}

Yeni bağlanan subscriber geçmişte publish edilmiş mesajları almaz.

Yalnızca bağlantı kurulduktan sonra publish edilen mesajlar alınır.

\subsection{TRANSIENT\_LOCAL}

Publisher son mesajı saklayabilir.

Yeni bağlanan subscriber, yeni bir mesaj beklemek zorunda kalmadan
publisher tarafından saklanan son değeri alabilir.

Bu davranış özellikle state machine, lamba ve sinyal durumlarının
aktarılmasında kullanılmaktadır.

% ================================================================
\section{Başlangıç Frekansları}
% ================================================================

Bu dokümanda tanımlanan frekanslar sistemin ilk entegrasyon aşaması için
başlangıç değerleridir.

Belirlenen temel frekanslar:

\begin{verbatim}
cmd_vel          : 20 Hz

autonomy_enable  : 5 Hz

motion_enable    : 10 Hz

vehicle speed    : 20 Hz

steering angle   : 20 Hz

IMU              : 50 Hz

distance sensors : 20 Hz

state machine    : On Change

signals/lights   : On Change
\end{verbatim}

Sistem testleri sırasında ihtiyaç görülmesi halinde bu değerler
güncellenebilir.

Özellikle araç hızı ve direksiyon açısı için başlangıçta 20 Hz
kullanılması planlanmıştır.

Daha yüksek örnekleme veya kontrol frekansına ihtiyaç duyulması durumunda
bu değerler daha sonra artırılabilir.

% ================================================================
\section{Sonuç}
% ================================================================

Bu dokümanda tanımlanan topic isimleri, mesaj tipleri, fiziksel birimler,
yayın frekansları ve QoS profilleri ROS 2 tarafındaki yazılım arayüz
sözleşmesini oluşturmaktadır.

ROS 2 tarafındaki:

\begin{itemize}

    \item otonom sürüş,

    \item kontrol,

    \item sensör işleme,

    \item durum izleme,

    \item kayıt,

    \item görselleştirme

\end{itemize}

düğümlerinin bu arayüz tanımlarına uygun şekilde geliştirilmesi
gerekmektedir.

Bridge, ROS 2 tarafı ile gömülü kontrol sistemi arasındaki arayüzü
sağlayacak ve ROS 2 düğümlerinin araç üzerindeki haberleşme altyapısının
detaylarından bağımsız olarak çalışmasını sağlayacaktır.

% ================================================================
\end{document}
